\section*{Capítulo \ref{ch:g6:solutions-and-solubility} --- Soluções e solubilidade}

\begin{solution}{exo:g6:solutions-and-solubility:1}
\omterm{def:g6:solutions-and-solubility:solute}{Soluto}: o açúcar. \omterm{def:g6:solutions-and-solubility:solute}{Solvente}: a água. Sim, o \omterm{def:g6:solutions-and-solubility:solute}{solvente} é a água: é uma
\omterm{def:g6:solutions-and-solubility:solute}{solução aquosa}.
\end{solution}

\begin{solution}{exo:g6:solutions-and-solubility:2}
Uma solução que não consegue \omterm{def:g2:mixing-and-dissolving:dissolve}{dissolver} mais nada do seu \omterm{def:g6:solutions-and-solubility:solute}{soluto}. O \omterm{def:g6:solutions-and-solubility:solute}{soluto}
acrescentado a ela fica no fundo, sem se \omterm{def:g2:mixing-and-dissolving:dissolve}{dissolver}.
\end{solution}

\begin{solution}{exo:g6:solutions-and-solubility:3}
$250 + 20 = \qty{270}{g}$.
\end{solution}

\begin{solution}{exo:g6:solutions-and-solubility:4}
Etanol: \omterm{def:g6:solutions-and-solubility:miscible}{miscível}. Óleo de cozinha: \omterm{def:g6:solutions-and-solubility:miscible}{imiscível}. Vinagre: \omterm{def:g6:solutions-and-solubility:miscible}{miscível} (ele é
principalmente água).
\end{solution}

\begin{solution}{exo:g6:solutions-and-solubility:5}
Propanona (acetona). GHS02: altamente inflamável; GHS07: irrita os
olhos, e o vapor pode provocar sonolência ou tontura.
\end{solution}

\begin{solution}{exo:g6:solutions-and-solubility:6}
$360 \times \frac{500}{1000} = \qty{180}{g}$.
\end{solution}

\begin{solution}{exo:g6:solutions-and-solubility:7}
\qty{200}{g} de água dissolvem no máximo $360 \times \frac{200}{1000} =
\qty{72}{g}$. Não: $80 - 72 = \qty{8}{g}$ ficam no fundo.
\end{solution}

\begin{solution}{exo:g6:solutions-and-solubility:8}
A solução pesa $225 + 25 = \qty{250}{g}$; $\frac{25}{250} =
\qty{10}{\percent}$.
\end{solution}

\begin{solution}{exo:g6:solutions-and-solubility:9}
No primeiro béquer, onde todo o sal se dissolveu e a solução ainda não
está saturada.
\end{solution}

\begin{solution}{exo:g6:solutions-and-solubility:10}
O dióxido de carbono foi dissolvido sob pressão na garrafa fechada. Ao
abri-la, a pressão diminui: a água consegue conter menos gás, e o gás em
excesso sai em forma de bolhas.
\end{solution}

\begin{solution}{exo:g6:solutions-and-solubility:11}
$2000 \div 0.04 = 50\,000$ vezes mais.
\end{solution}

\begin{solution}{exo:g6:solutions-and-solubility:12}
Um gás \omterm{def:g2:mixing-and-dissolving:dissolve}{se dissolve} menos na água morna: o lago aquecido contém menos
oxigênio dissolvido, e falta oxigênio para os peixes.
\end{solution}

\begin{solution}{pb:g6:solutions-and-solubility:1}
\textbf{1.} \omterm{def:g6:solutions-and-solubility:solute}{Soluto}: o sal. \omterm{def:g6:solutions-and-solubility:solute}{Solvente}: a água.

\textbf{2.} A água contém todo o sal que consegue conter: mais nenhum sal
pode se \omterm{def:g2:mixing-and-dissolving:dissolve}{dissolver} nela.

\textbf{3.} O sal não dissolvido no fundo prova que a salmoura está
saturada: ela contém o máximo de sal possível.

\textbf{4.} $360 \times 20 = \qty{7200}{g} = \qty{7.2}{kg}$.

\textbf{5.} $20 + 7.2 = \qty{27.2}{kg}$.

\textbf{6.} $\frac{7.2}{27.2} \approx 0.26$: cerca de \qty{26}{\percent}.

\textbf{7.} Não: a água poderia \omterm{def:g2:mixing-and-dissolving:dissolve}{dissolver} \qty{7.2}{kg}. Ela ainda
poderia \omterm{def:g2:mixing-and-dissolving:dissolve}{dissolver} $7.2 - 5 = \qty{2.2}{kg}$.

\textbf{8.} $20 - 5 = \qty{15}{kg}$ de água (\qty{15}{L}).

\textbf{9.} $360 \times 15 = \qty{5400}{g} = \qty{5.4}{kg}$.

\textbf{10.} Ele sai da solução em forma de cristais sólidos: ele
cristaliza e desce para o fundo.

\textbf{11.} Aparecem cristais brancos novos, que se acumulam no fundo do
tanque.

\textbf{12.} $7.2 - 5.4 = \qty{1.8}{kg}$ de sal cristalizaram.
\end{solution}
